% ============================================================================
% THE 512-BIT CHORD REPRESENTATION
% ============================================================================
\subsection{Deterministic Prime-Basis Identity}
\label{sec:base_chord}
\label{sec:chord_representation}

Let $D = 512$ denote the chord dimensionality and let $\mathcal{B} = \{0,
\ldots, 255\}$ be the byte alphabet.  Every byte value $b \in \mathcal{B}$
is mapped to a deterministic binary vector $\mathbf{c}(b) \in \{0,1\}^{D}$
by activating exactly five bits at positions given by coprime spreading:

\begin{equation}
  \mathcal{P}(b) = \bigl\{7b \bmod D,\; 13b \bmod D,\; 31b \bmod D,\;
                          61b \bmod D,\; 127b \bmod D\bigr\}
  \label{eq:base_chord}
\end{equation}

\begin{equation}
  c_i(b) =
  \begin{cases}
    1 & \text{if } i \in \mathcal{P}(b), \\
    0 & \text{otherwise.}
  \end{cases}
  \label{eq:indicator}
\end{equation}

\noindent Because the multipliers $\{7, 13, 31, 61, 127\}$ are pairwise
coprime and $D = 2^9$, distinct byte values receive distinct activation
patterns (verified by exhaustive enumeration over all $256$ values).  The
resulting chord is extremely sparse: $\|\mathbf{c}(b)\|_1 = 5$ out of
$D = 512$ bits, yielding a sparsity ratio of $5/512 \approx 0.98\%$.

\paragraph{Why Coprime Multipliers.}
The choice of coprime multipliers is not arbitrary.  Multiplying by a prime
$p$ and reducing modulo a power of two $D = 2^k$ generates a permutation
of the residues $\{0, \ldots, D-1\}$ if and only if $\gcd(p, D) = 1$.
Since all five multipliers are odd primes and $D$ is a power of two, each
multiplier independently generates a distinct full-period permutation,
maximizing the minimum Hamming distance between any two byte chords.

\subsection{Algebraic Operations}
\label{sec:chord_algebra}

The Chord supports a closed algebra of compositional operations.  Each
executes in $O(D/w)$ time where $w$ is the SIMD word width (64 for
\texttt{uint64}, yielding $D/w = 8$ operations per chord):

\begin{align}
  \textbf{ChordOR:}\quad
    (\mathbf{a} \lor \mathbf{b})_i
    &= a_i \lor b_i
    &&\text{(Superposition / LCM)}
    \label{eq:chord_or} \\[4pt]
  \textbf{ChordGCD:}\quad
    (\mathbf{a} \land \mathbf{b})_i
    &= a_i \land b_i
    &&\text{(Shared factors / Intersection)}
    \label{eq:chord_gcd} \\[4pt]
  \textbf{ChordHole:}\quad
    H(\mathbf{t}, \mathbf{e})_i
    &= t_i \land \lnot e_i
    &&\text{(Structural Vacuum)}
    \label{eq:chord_hole} \\[4pt]
  \textbf{Similarity:}\quad
    \text{sim}(\mathbf{a}, \mathbf{b})
    &= \sum_{i=0}^{D-1} a_i \cdot b_i
    = \texttt{popcount}(\mathbf{a} \land \mathbf{b})
    &&\text{(Discrete dot product)}
    \label{eq:chord_sim}
\end{align}

\paragraph{Equivalence to Inner Product.}
For binary vectors $\mathbf{a}, \mathbf{b} \in \{0,1\}^D$, the real-valued
inner product $\langle \mathbf{a}, \mathbf{b} \rangle$ reduces exactly to
the popcount of their bitwise AND:
\begin{equation}
  \langle \mathbf{a}, \mathbf{b} \rangle
  = \sum_{i=0}^{D-1} a_i \, b_i
  = \bigl| \{i : a_i = 1 \land b_i = 1\} \bigr|
  = \texttt{popcount}(\mathbf{a} \mathbin{\&} \mathbf{b})
  \label{eq:inner_product}
\end{equation}
This identity is the mechanism by which single-cycle integer instructions
execute semantic similarity computation.

\subsection{Positional Binding via Circular Shift}
\label{sec:roll_left}

To encode sequential position without destroying content, we define the
\emph{RollLeft} operator as a circular permutation over bit indices:

\begin{equation}
  \text{RollLeft}(\mathbf{c},\, s)_i = c_{(i - s) \bmod D}
  \label{eq:roll_left}
\end{equation}

Because circular shifts are invertible and non-commutative with
superposition, positional binding preserves all algebraic properties while
ensuring that re-ordered sequences produce distinct chords:

\begin{equation}
  \text{RollLeft}(\mathbf{c}_d, 0) \lor \text{RollLeft}(\mathbf{c}_o, 1)
  \lor \text{RollLeft}(\mathbf{c}_g, 2)
  \;\neq\;
  \text{RollLeft}(\mathbf{c}_g, 0) \lor \text{RollLeft}(\mathbf{c}_o, 1)
  \lor \text{RollLeft}(\mathbf{c}_d, 2)
  \label{eq:dog_god}
\end{equation}

\noindent That is, $\texttt{``dog''} \neq \texttt{``god''}$ in the chord
representation.

\subsection{SimHash Structural Binning}
\label{sec:chord_bin}

The \emph{ChordBin} operator maps any chord $\mathbf{c}$ to an 8-bit
structural bin $b \in \{0, \ldots, 255\}$ via a SimHash-style random
hyperplane projection.  Given $k = 8$ fixed seed values
$\{s_1, \ldots, s_8\}$, the accumulator for plane $j$ sums signed
contributions from each active bit:

\begin{equation}
  \text{acc}_j = \sum_{i : c_i = 1}
  \begin{cases}
    +1 & \text{if } h(i, s_j) \geq 2^{63}, \\
    -1 & \text{otherwise,}
  \end{cases}
  \qquad
  h(i, s_j) = i \cdot s_j + (i \ll (j+1))
  \label{eq:simhash}
\end{equation}

The bin is then the 8-bit integer whose $j$-th bit is $[\text{acc}_j \geq 0]$.
Chords with high Hamming similarity map to the same or nearby bins with
high probability, enabling the EigenMode phase table
(\Cref{sec:eigenmode_cooccurrence}) to use a fixed-size $256 \times 256$
co-occurrence matrix without per-chord storage.

\subsection{FlatChord: GPU-Optimized Dense Representation}
\label{sec:flat_chord}

For GPU dispatch, the sparse 512-bit Chord is converted to a \emph{FlatChord}:
a dense array of at most $D$ active prime indices $\{p_1, \ldots, p_k\}$
where $k = \|\mathbf{c}\|_1$.  This eliminates bit-twiddling thread
divergence in SIMT architectures: every thread iterates over exactly
$k$ entries rather than scanning all $D/64 = 8$ blocks.

Batched flattening across $N$ chords is parallelized via a pre-warmed,
auto-scaling worker pool, yielding amortized $O(k)$ per chord with
near-zero goroutine creation overhead.
